% Generated draft: requires mathematical and source review.
\documentclass[11pt]{article}
\usepackage[margin=25mm]{geometry}
\usepackage{fontspec}
\setmainfont{Latin Modern Roman}
\usepackage{amsmath,amssymb,mathtools,unicode-math}
\setmathfont{Latin Modern Math}
\usepackage{enumitem,needspace,xurl,hyperref}
\setlength{\parindent}{0pt}
\setlength{\parskip}{5pt}
\setlength{\emergencystretch}{3em}
\newcommand{\N}{\mathbb{N}}
\newcommand{\Z}{\mathbb{Z}}
\newcommand{\Q}{\mathbb{Q}}
\newcommand{\R}{\mathbb{R}}
\newcommand{\C}{\mathbb{C}}
\newcommand{\F}{\mathbb{F}}
\newcommand{\A}{\mathbb{A}}
\newcommand{\PP}{\mathbb{P}}
\newcommand{\E}{\mathbb{E}}
\newcommand{\eps}{\varepsilon}
\newcommand{\dd}{\,\mathrm{d}}
\begin{document}
\section*{UnsolvedMath ID: 30006001}
\textbf{Source problem code:} catalogue entry
\subsection{Definitions and mathematical statement}
Let \(\Omega\subset\R^n\) be a Lipschitz domain: an open connected set whose boundary is locally, after a rigid change of coordinates, the graph of a Lipschitz function. No smallness assumption is imposed on the Lipschitz constants. The supplied statement imposes neither boundedness of \(\Omega\) nor an explicit restriction on the ambient dimension \(n\). Let \(\Sigma\subset\partial\Omega\) be relatively open, and let \(u\in C(\overline\Omega)\) satisfy
\[
\Delta u=0\quad\text{in }\Omega,
\qquad u|_{\Sigma}=0.
\]
Here \(\mathcal H^{n-1}\) denotes \((n-1)\)-dimensional Hausdorff measure.

The intended boundary uniqueness question is whether
\[
\mathcal H^{n-1}\bigl(\{x\in\Sigma:\nabla u(x)=0\}\bigr)>0
\quad\Longrightarrow\quad
u\equiv0\text{ in }\Omega.
\]
\textbf{Unresolved boundary interpretation.} The displayed boundary value of \(\nabla u\) is not defined by the stated hypotheses alone. Continuity of \(u\) on \(\overline\Omega\) and harmonicity in \(\Omega\) do not specify a boundary gradient trace. To make the intended question precise after choosing a trace convention \(T\), define
\[
E_T=\{x\in\Sigma:T(\nabla u)(x)\text{ exists and equals }0\}.
\]
The corresponding formulation asks, for every \(\Omega,\Sigma,u\) satisfying the hypotheses above, whether
\[
\mathcal H^{n-1}(E_T)>0
\quad\Longrightarrow\quad
u\equiv0\text{ in }\Omega.
\]
The convention \(T\), including its existence and measurability requirements, must be recovered from the cited report before this becomes a fully specified problem. Possible interpretations include nontangential boundary limits or a suitably justified weak trace; none is selected by the supplied record. No additional differentiability or Sobolev assumption is asserted here.

The supplied original passage asks whether an earlier statement remains valid when the smallness assumption on the Lipschitz constant is removed, but that earlier statement is not reproduced. Thus the reconstructed formulation and its trace hypotheses require human comparison with the report.
\subsection{Short English statement}
Must a harmonic function, continuous up to the boundary of an arbitrary Lipschitz domain and zero on a relatively open boundary patch, vanish identically if its boundary gradient is zero on a set of positive surface measure? The intended boundary gradient trace is unspecified. The supplied curation describes this as removing the small Lipschitz slope assumption from a theorem attributed to Tolsa.
\subsection{Sources}
\begin{itemize}
\item Supplied local UnsolvedMath registry record, permanent ID \(30006001\), title ``Boundary Uniqueness on Lipschitz Domains,'' source problem code \texttt{catalogue entry}, in the Oberwolfach Reports Open Problems collection. The record gives proposed year 2024 and names no proposer.
\item Supplied report citation: \emph{Calculus of Variations (2025)}, Oberwolfach Reports, DOI \texttt{10.4171/owr/2024/37}, pp.\ 2113--2178. \url{https://doi.org/10.4171/owr/2024/37}. The supplied title year and DOI/report year differ; both are retained without reconciliation.
\item Supplied literature source: Oberwolfach Report 37/2024. \url{https://ems.press/content/serial-article-files/50045}.
\item Supplied literature source: Oberwolfach Report 37/2022. \url{https://ems.press/content/serial-article-files/46973}.
\item Literature assessment dated \texttt{2026-08-22}: the supplied curation labels the problem open and states that Tolsa's boundary uniqueness theorem requires sufficiently small Lipschitz slope, with removal of that assumption remaining open. This is supplied dated curation, not an independently verified literature assessment. The record itself requests independent verification of time-sensitive status and cited claims.
\item The registry labels the statement \texttt{corrected\_verified} and reports that it was reconstructed and checked against the official Oberwolfach report. These are supplied verification claims; no independent source inspection is asserted in this draft.
\end{itemize}
\end{document}
